\section{考察}
\subsection{客観評価と主観評価の乖離}
報告者が担当したトロンボーンおよびカスタネットは，客観評価においては周波数誤差$\pm1.0$\,\%以内という
基準を満たした一方で，主観評価では作成音の平均評価値がいずれも基準の4.0に達しなかった．この乖離に
ついて考察する．

ここで，客観評価と主観評価では測定条件が異なる点に注意する必要がある．客観評価は音源単体で取得した
波形に基づくものであり，音源の周波数特性そのものを評価している．これに対し主観評価は，人形の動作を
含むシステム全体の演奏を提示して行ったものである．したがって，主観評価には音源そのものの音色だけで
なく，演奏環境全体の要因が反映される．

主観評価が伸びなかった要因としては，合成音の音量が実際の演奏環境において十分でなかったこと，および
他楽器と比べ，オクターブを一つ下で設定しているため，音色が低く聞こえ，音量が小さく感じられたことが考えられる．
特にヴァイオリンは高音域で音色的にも目立つ楽器なので，トロンボーンやカスタネットの音が埋もれてしまい，主観評価において低い評価となったと考えられる．
これに関してはピアノでも同様の傾向が見られたため，ヴァイオリンの音が他楽器の音が埋もれる原因となった可能性が高いと思われる．

\subsection{トロンボーンの音色成立の要因}
トロンボーンの音源では，波形テーブルに設定した倍音構成が，実際のトロンボーンの倍音構造を正確に
反映したものではなかった．にもかかわらず，合成音は電子的なブザー音ではなく，
管楽器に近い音色として成立し，主観評価においても音質は3.97と比較的高い評価を得た．この要因について
考察する．

本音源では，オシレータが生成する波形に対し，ローパスフィルタ，ディレイ，および音量エンベロープを
直列に適用している．ローパスフィルタは高次の倍音成分を抑制するため，波形テーブルの倍音構成が
正確でなくとも，過剰な高域成分が除去され，硬い電子音的な質感が緩和される．また，音量エンベロープに
よる緩やかな立ち上がりと減衰は，管楽器特有の発音の印象を与える．これらのことから，本音源の音色は，
波形テーブルの倍音構成そのものよりも，フィルタおよびエンベロープによる後段の処理に強く依存して
決定されていたと考えられる．この結果は，オシレータの倍音設計が不完全な場合でも，フィルタと
エンベロープの設計によって楽器らしい音色を得られる可能性を示している．

\subsection{カスタネットの打撃音合成の有効性}
カスタネットの音源では，オシレータではなくノイズを音源とし，これを帯域通過フィルタと音量エンベロープに
通す構成を採用した．評価の結果，合成音は立ち上がり直後に振幅が最大となり，約20\,msで減衰する打撃音
特有の波形を示した．これは，音程を持たない非調和的な打撃音という，カスタネットの音響的特徴を
再現できたことを示している．オシレータを用いる合成では明確な基本周波数を持つ音となり，打楽器の
質感を再現することが難しいのに対し，ノイズを音源とし帯域通過フィルタで周波数帯を限定する手法は，
打楽器音の合成に有効であったと考えられる．

\subsection{改善案}
以上の評価と考察を踏まえ，今後の改善点を述べる．第一に，主観評価が伸びなかった要因である音量の
不足については，音源の出力レベルを見直すとともに，Processing側でゲインを調整することでの改善が期待される．
第二に，トロンボーンの音色については，実際の楽器の倍音構造をピーク検出によって正確に抽出し，
波形テーブルに反映することで，より実物に近い音色が得られると考えられる．これらの改善により，
客観評価に加えて主観評価においても基準を満たす音源の実現が期待される．